\documentclass{amsart}
% For dealing with references we use the comment environment
\usepackage{verbatim}
\newenvironment{reference}{\comment}{\endcomment}
%\newenvironment{reference}{}{}
\newenvironment{slogan}{\comment}{\endcomment}
\newenvironment{history}{\comment}{\endcomment}

% For commutative diagrams we use Xy-pic
\usepackage[all]{xy}

% We use 2cell for 2-commutative diagrams.
\xyoption{2cell}
\UseAllTwocells

% We use multicol for the list of chapters between chapters
\usepackage{multicol}

% This is generally recommended for better output
\usepackage{lmodern}
\usepackage[T1]{fontenc}

% For cross-file-references

% Package for hypertext links:
\usepackage{hyperref}

% For any local file, say "hello.tex" you want to link to please

% Theorem environments.
%
\theoremstyle{plain}
\newtheorem{theorem}[subsection]{Theorem}
\newtheorem{proposition}[subsection]{Proposition}
\newtheorem{lemma}[subsection]{Lemma}

\theoremstyle{definition}
\newtheorem{definition}[subsection]{Definition}
\newtheorem{example}[subsection]{Example}
\newtheorem{exercise}[subsection]{Exercise}
\newtheorem{situation}[subsection]{Situation}

\theoremstyle{remark}
\newtheorem{remark}[subsection]{Remark}
\newtheorem{remarks}[subsection]{Remarks}

\numberwithin{equation}{subsection}

% Macros
%
\def\lim{\mathop{\mathrm{lim}}\nolimits}
\def\colim{\mathop{\mathrm{colim}}\nolimits}
\def\Spec{\mathop{\mathrm{Spec}}}
\def\Hom{\mathop{\mathrm{Hom}}\nolimits}
\def\Ext{\mathop{\mathrm{Ext}}\nolimits}
\def\SheafHom{\mathop{\mathcal{H}\!\mathit{om}}\nolimits}
\def\SheafExt{\mathop{\mathcal{E}\!\mathit{xt}}\nolimits}
\def\Sch{\mathit{Sch}}
\def\Mor{\mathop{\mathrm{Mor}}\nolimits}
\def\Ob{\mathop{\mathrm{Ob}}\nolimits}
\def\Sh{\mathop{\mathit{Sh}}\nolimits}
\def\NL{\mathop{N\!L}\nolimits}
\def\CH{\mathop{\mathrm{CH}}\nolimits}
\def\proetale{{pro\text{-}\acute{e}tale}}
\def\etale{{\acute{e}tale}}
\def\QCoh{\mathit{QCoh}}
\def\Ker{\mathop{\mathrm{Ker}}}
\def\Im{\mathop{\mathrm{Im}}}
\def\Coker{\mathop{\mathrm{Coker}}}
\def\Coim{\mathop{\mathrm{Coim}}}

% Boxtimes
%
\DeclareMathSymbol{\boxtimes}{\mathbin}{AMSa}{"02}

%
% Macros for moduli stacks/spaces
%
\def\QCohstack{\mathcal{QC}\!\mathit{oh}}
\def\Cohstack{\mathcal{C}\!\mathit{oh}}
\def\Spacesstack{\mathcal{S}\!\mathit{paces}}
\def\Quotfunctor{\mathrm{Quot}}
\def\Hilbfunctor{\mathrm{Hilb}}
\def\Curvesstack{\mathcal{C}\!\mathit{urves}}
\def\Polarizedstack{\mathcal{P}\!\mathit{olarized}}
\def\Complexesstack{\mathcal{C}\!\mathit{omplexes}}
% \Pic is the operator that assigns to X its picard group, usage \Pic(X)
% \Picardstack_{X/B} denotes the Picard stack of X over B
% \Picardfunctor_{X/B} denotes the Picard functor of X over B
\def\Pic{\mathop{\mathrm{Pic}}\nolimits}
\def\Picardstack{\mathcal{P}\!\mathit{ic}}
\def\Picardfunctor{\mathrm{Pic}}
\def\Deformationcategory{\mathcal{D}\!\mathit{ef}}

\setlength{\emergencystretch}{3em}
\begin{document}
\title[Stacks Project, Tag 07PM]{347 --- Regularity of flat complete local maps with geometrically regular closed fibre}
\maketitle
\noindent Copyright (C) 2005--2025 Johan de Jong. Stacks Project authors. Source: \href{https://stacks.math.columbia.edu/tag/07PM}{Tag 07PM}.

\noindent Editorial difficulty note (not author text). Difficulty H2: The author must prove geometric regularity at every prime, not just at the closed fibre. It reduces to a power-series base and inducts through purely inseparable extensions, extending derivations to detect each degree-p equation and preserve regularity..

\noindent Editorial provenance note (not author text). History: excerpt from revision \texttt{a0444\allowbreak 6e57e\allowbreak c1fbc\allowbreak 252a8\allowbreak 71afc\allowbreak ec775\allowbreak 2fb28\allowbreak 07b14}, more-algebra.tex, lines 12027--12116. Reference presentation was adapted; original-author context and core support blocks were retained or added. The mathematical wording is unchanged. Licensed under GNU FDL 1.2 or later; no invariant sections or cover texts. See ../licenses/COPYING. Context blocks reproduce author definitions and supporting results; they are not additional counted problems.

\begin{proposition}
\label{proposition-fs-regular}
Let $A \to B$ be a local homomorphism of Noetherian complete local rings.
Let $k$ be the residue field of $A$ and $\overline{B} = B \otimes_A k$
the special fibre.
The following are equivalent
\begin{enumerate}
\item $A \to B$ is regular,
\item $A \to B$ is flat and $\overline{B}$ is geometrically regular
over $k$,
\item $A \to B$ is flat and $k \to \overline{B}$ is formally smooth
in the $\mathfrak m_{\overline{B}}$-adic topology, and
\item $A \to B$ is formally smooth in the $\mathfrak m_B$-adic
topology.
\end{enumerate}
\end{proposition}

\begin{proof}
We have seen the equivalence of (2), (3), and (4) in
Proposition \href{https://stacks.math.columbia.edu/tag/07NQ}{proposition fs flat fibre fs (Tag 07NQ)}.
It is clear that (1) implies (2).
Thus we assume the equivalent conditions (2), (3), and (4) hold
and we prove (1).

\medskip\noindent
Let $\mathfrak p$ be a prime of $A$. We will show that
$B \otimes_A \kappa(\mathfrak p)$ is geometrically regular
over $\kappa(\mathfrak p)$.
By Lemma \href{https://stacks.math.columbia.edu/tag/07EG}{lemma base change fs (Tag 07EG)}
we may replace $A$ by $A/\mathfrak p$ and $B$ by $B/\mathfrak pB$.
Thus we may assume that $A$ is a domain and that $\mathfrak p = (0)$.

\medskip\noindent
Choose $A_0 \subset A$ as in Algebra, Lemma
\href{https://stacks.math.columbia.edu/tag/032D}{lemma complete local Noetherian domain finite over regular (Tag 032D)}.
We will use all the properties stated in that lemma without further mention.
As $A_0 \to A$ induces an isomorphism on residue fields, and as
$B/\mathfrak m_A B$ is geometrically regular over $A/\mathfrak m_A$
we can find a diagram
$$
\xymatrix{
C \ar[r] & B \\
A_0 \ar[r] \ar[u] & A \ar[u]
}
$$
with $A_0 \to C$ formally smooth in the $\mathfrak m_C$-adic topology
such that $B = C \otimes_{A_0} A$, see Remark \href{https://stacks.math.columbia.edu/tag/07NS}{remark what does it mean (Tag 07NS)}.
(Completion in the tensor product is not needed as $A_0 \to A$ is
finite, see Algebra, Lemma \href{https://stacks.math.columbia.edu/tag/00MA}{lemma completion tensor (Tag 00MA)}.)
Hence it suffices to show that $C \otimes_{A_0} K_0$
is a geometrically regular algebra over the fraction field $K_0$ of $A_0$.

\medskip\noindent
The upshot of the preceding paragraph is that we may assume that
$A = k[[x_1, \ldots, x_n]]$ where $k$ is a field or
$A = \Lambda[[x_1, \ldots, x_n]]$ where $\Lambda$ is a Cohen ring.
In this case $B$ is a regular ring, see
Algebra, Lemma \href{https://stacks.math.columbia.edu/tag/031E}{lemma flat over regular with regular fibre (Tag 031E)}.
Hence $B \otimes_A K$ is a regular ring too (where $K$ is the
fraction field of $A$) and we win
if the characteristic of $K$ is zero.

\medskip\noindent
Thus we are left with the case where $A = k[[x_1, \ldots, x_n]]$
and $k$ is a field of characteristic $p > 0$.
Let $L/K$ be a finite purely inseparable field extension.
We will show by induction on $[L : K]$ that $B \otimes_A L$
is regular. The base case is $L = K$ which we've seen above.
Let $K \subset M \subset L$ be a subfield such that
$L$ is a degree $p$ extension of $M$ obtained by adjoining a $p$th root
of an element $f \in M$. Let $A'$ be a finite $A$-subalgebra
of $M$ with fraction field $M$. Clearing denominators, we may and do assume
$f \in A'$. Set $A'' = A'[z]/(z^p -f)$ and note that $A' \subset A''$
is finite and that the fraction field of $A''$ is $L$.
By induction we know that $B \otimes_A M$ ring is regular.
We have
$$
B \otimes_A L = B \otimes_A M[z]/(z^p - f)
$$
By Lemma \href{https://stacks.math.columbia.edu/tag/07PH}{lemma find D (Tag 07PH)} we know there exists a derivation
$D : A' \to A'$ such that $D(f) \not = 0$. As $A' \to B \otimes_A A'$
is formally smooth in the $\mathfrak m$-adic topology by
Lemma \href{https://stacks.math.columbia.edu/tag/07EH}{lemma descent fs (Tag 07EH)}
we can use
Lemma \href{https://stacks.math.columbia.edu/tag/07PL}{lemma lift derivation through fs (Tag 07PL)}
to extend $D$ to a derivation $D' : B \otimes_A A' \to B \otimes_A A'$.
Note that $D'(f) = D(f)$ is a unit in $B \otimes_A M$ as $D(f)$
is not zero in $A' \subset M$. Hence $B \otimes_A L$ is regular by
Lemma \href{https://stacks.math.columbia.edu/tag/07PG}{lemma degree p extension regular (Tag 07PG)} and we win.
\end{proof}
\end{document}
